%% ICML-style single-blind submission format.
%% Self-contained: compiles with any standard LaTeX install (no icml20xx.sty needed).
%% To use the official style instead, see README.md in this directory.
\documentclass[twocolumn,10pt]{article}

\usepackage[margin=0.75in]{geometry}
\usepackage{times}
\usepackage{amsmath,amssymb}
\usepackage{booktabs}
\usepackage{graphicx}
\usepackage{microtype}
\usepackage[hidelinks]{hyperref}
\usepackage{url}
\usepackage{caption}
\usepackage{natbib}
\usepackage{titlesec}

\setlength{\columnsep}{0.25in}
\captionsetup{font=small,labelfont=bf,skip=4pt}
\titleformat{\section}{\normalfont\large\bfseries}{\thesection}{0.5em}{}
\titleformat{\subsection}{\normalfont\normalsize\bfseries}{\thesubsection}{0.5em}{}
\renewcommand{\abstractname}{\normalsize Abstract}
\bibliographystyle{plainnat}

\title{\bf When the Baseline Has No Parameters:\\
Evaluation Pitfalls in Single-Event Industrial Anomaly Detection}

\author{Anonymous Submission}
\date{}

\begin{document}
\twocolumn[
\begin{@twocolumnfalse}
\maketitle
\begin{center}\begin{minipage}{0.86\textwidth}
\begin{center}\textbf{Abstract}\end{center}
\small\noindent
We report a methodological audit of a widely distributed industrial
predictive-maintenance benchmark: vibration and current time series from the
hydraulic pump of a fine-blanking press. The benchmark is representative of a
large class of public manufacturing datasets---one labelled fault episode,
one unlabelled normal recording, acquired on different days. We make four
observations that we believe generalise beyond this dataset. \emph{(i)} The
acquisition is burst-sampled: $58\%$ of wall-clock time is gaps, and naive
sliding windows span those gaps, inflating the window count by $2\times$ and
manufacturing detectable amplitude discontinuities; gap-aware windowing lowers
the reference $F_1$ from $0.857$ to $0.818$, and the entire loss is in recall.
\emph{(ii)} Window-level $F_1$ is saturated: a zero-parameter range rule
attains $F_1=0.857$ against $0.748$ for the reference LSTM-autoencoder, so
reporting $F_1$ alone cannot distinguish learned models from trivial ones.
\emph{(iii)} Three commonly used operational metrics degenerate when the
positive class is a single event---detection lead time is \emph{unattainable by
construction}, event recall is $1.0$ for every detector, and a false-alarm rate
of $1$/hour is unresolvable at the available $0.139$\,h of normal observation
(Poisson $95\%$ upper bound $21.6$/h at zero alarms). \emph{(iv)} The current
channel separates the two recordings perfectly ($F_1=1.000$ from two features
alone), but its mains carrier occupies $95.9\%$ of normal AC power versus
$10.4\%$ in the fault recording---physically impossible for a running motor,
and therefore an acquisition artefact rather than a fault signature. We adopt a
pre-declared exclusion rule, quantify its cost at $\Delta F_1 = 0.072$, and show
that a vibration-only detector still reaches $F_1=0.928$ (95\% CI
$0.852$--$0.984$, zero false positives, zero learned parameters). Finally, we
localise $91.1\%$ of all false alarms to a single reversible $300$\,s low-load
excursion inside the \emph{normal} recording, which we argue is the genuinely
informative finding the benchmark affords.
\end{minipage}\end{center}
\vspace{1.5em}
\end{@twocolumnfalse}
]

\section{Introduction}

Public benchmarks for industrial anomaly detection are typically small, weakly
labelled, and assembled from a handful of recording sessions. A recurring
pattern---visible in the dataset we study and in many of its peers---is that
the ``normal'' and ``anomalous'' classes correspond to \emph{different
recording sessions}. Any difference between sessions, whether mechanical or
instrumental, is then perfectly confounded with the label.

This is not a hypothetical concern. \citet{wu2021flawed} showed that several
widely used time-series anomaly benchmarks are solvable by one-line rules;
\citet{kim2022rigorous} showed that the point-adjustment protocol common in the
literature inflates $F_1$ to the point where random scores outperform published
models. Our contribution is complementary and concrete: we take a single
benchmark that is actively used in a national competition, and we trace, with
measured numbers, exactly which evaluation decisions change the conclusion and
by how much.

We organise the paper around the decisions an analyst must make, in the order
they arise: how to window burst-sampled data (\S\ref{sec:data}), which metrics
remain identifiable when the positive class is one event (\S\ref{sec:metrics}),
how to decide whether a highly discriminative channel is signal or artefact
(\S\ref{sec:confound}), and only then which detector to use
(\S\ref{sec:method}--\S\ref{sec:results}). Our view is that the first three
decisions dominate the fourth, and that reporting practice should reflect that
ordering.

\paragraph{Contributions.}
\begin{itemize}\itemsep1pt
\item We identify a burst-sampling acquisition structure and show that
gap-unaware windowing discards the distinction between adjacent and
non-adjacent samples, inflating recall through artefacts (\S\ref{sec:data}).
\item We show window $F_1$ on this benchmark is saturated by a zero-parameter
rule, and propose a two-tier baseline protocol that preserves both an honest
improvement margin and a strict reference (\S\ref{sec:metrics}).
\item We give a falsifiable, pre-declared exclusion rule for session-confounded
channels and report the exclusion cost rather than hiding it
(\S\ref{sec:confound}).
\item We localise the dominant false-alarm mode to a reversible operating-mode
excursion inside the normal class, and show it is the mechanism behind the
apparent robustness of amplitude features (\S\ref{sec:fp}).
\end{itemize}

\section{Dataset and Acquisition Structure}
\label{sec:data}

\subsection{The benchmark}

The dataset provides two CSV files sampled at a nominal $10$\,Hz with three
analogue channels---upper motor vibration (\textsf{AI0}), lower motor vibration
(\textsf{AI1}), and motor current (\textsf{AI2})---plus a constant
\textsf{Equipment\_state} flag. The normal file contains $20{,}000$ rows
recorded on 2022-07-12 ($77$\,min wall clock); the fault file contains $600$
rows recorded on 2022-07-17 ($2.8$\,min). Every row of the normal file is
labelled $0$ and every row of the fault file is labelled $1$: the label is a
\emph{file} attribute, not a row attribute. There is therefore exactly one
positive event.

\subsection{Burst sampling}

Consecutive-sample intervals are not uniform. Taking a gap threshold of
$0.5$\,s, the normal file decomposes into $599$ contiguous runs and the fault
file into $21$. Run lengths are capped at exactly $50$ samples ($5$\,s) in both
files, with $33.7\%$ of normal runs attaining the cap. Gaps account for
$58.0\%$ of normal wall-clock time and $65.0\%$ of fault wall-clock time. The
inter-burst start interval is tightly distributed (median $7.96$\,s) while
burst length is not, which is the signature of a periodic publish-and-drop
transport layer rather than a process cycle.

This matters for windowing. With $L=20$ samples, only $9{,}978$ of the $19{,}981$
naive stride-1 windows lie entirely inside a run ($49.9\%$ survival); at $L=10$
survival is $74.4\%$ and at $L=5$ it is $88.3\%$. A naive window that straddles
a gap concatenates signal segments separated by up to $16.6$\,s and so contains
a discontinuity of purely instrumental origin.

Two consequences follow. First, nominal window duration is wrong: a naive
$20$-sample window spans a median $3.35$\,s (max $24.8$\,s) of wall clock, not
$2$\,s; gap-aware windows span exactly $1.9$\,s. Second---and this is the
quantitative point---removing straddling windows lowers the reference
zero-parameter detector from $F_1=0.857$ to $0.818$, and the drop is
\emph{entirely} in recall ($0.750\rightarrow0.692$) while precision and FPR stay
at $1.000$ and $0.000$. The discarded windows were therefore not informative
detections but gap-induced amplitude jumps.

\begin{table}[t]
\centering\small
\caption{Effect of windowing on the zero-parameter range rule. All rows use
identical splits and thresholds; only the windowing differs. The recall-only
loss indicates that gap-straddling windows were being detected for
instrumental, not mechanical, reasons.}
\label{tab:windowing}
\begin{tabular}{lcccc}
\toprule
Windowing & $N_{\text{norm}}/N_{\text{fault}}$ & P & R & $F_1$ \\
\midrule
naive, offset 100 & 4000 / 180 & 1.000 & 0.750 & 0.857 \\
naive, no reserve & 4880 / 480 & 1.000 & 0.704 & 0.826 \\
gap-aware, offset 100 & 2433 / 239 & 1.000 & 0.649 & 0.787 \\
\textbf{gap-aware} & \textbf{2489 / 276} & \textbf{1.000} & \textbf{0.692} & \textbf{0.818} \\
\bottomrule
\end{tabular}
\end{table}

\subsection{Spectral structure}

The reference analysis distributed with the dataset characterises normal
vibration as uncorrelated. Inspecting lag-1 autocorrelation alone supports
this ($r_1 = -0.011$ for \textsf{AI0}), but it is an artefact of the chosen lag.
Burst-averaged periodograms ($50$ samples, Hann window, $0.2$\,Hz resolution,
uniform expectation $3.8\%$ per bin) place $25.0\%$ of \textsf{AI0} power in a
single $3.6$\,Hz bin and $38.0\%$ of \textsf{AI1} power in a single $1.8$\,Hz
bin, in a stable $2{:}1$ ratio present in every burst. A Ljung--Box test
\citep{ljung1978measure} over lags $1$--$8$ rejects whiteness in $275$ of $276$
normal \textsf{AI0} bursts. Normal operation is thus a narrowband line
spectrum, and the fault condition is the \emph{collapse} of that line structure
into broadband low-frequency content.

We note that at $10$\,Hz the Nyquist limit is $5$\,Hz, whereas a gear-type
hydraulic pump driven by a four-pole induction motor has a gear-mesh frequency
of roughly $300$\,Hz. Bearing and gear fault frequencies are therefore not
recoverable from the vibration channels under any transform, which rules out
the envelope-spectrum family of methods regardless of modelling choice. We make
this explicit because the conclusion is insensitive to the pole-count
assumption: $2$-pole ($60$\,Hz) and $6$-pole ($20$\,Hz) shafts both exceed
Nyquist.

\section{Which Metrics Survive One Event?}
\label{sec:metrics}

\subsection{Degenerate operational metrics}

Operational anomaly-detection practice favours event-level metrics over
sample-level ones \citep{hundman2018detecting,garg2021evaluation}. On this
benchmark three such metrics are not merely noisy but \emph{unidentifiable}.

\paragraph{Detection lead time is unattainable.} Lead time is defined as
$t_0 - t_{\text{alarm}}$ where $t_0$ is fault onset. The fault file is labelled
positive from its first row, so $t_0$ can only be taken as that row and no
alarm can precede it. A window of $L=20$ requires $1.9$\,s to fill and a
$k=5$ debounce a further $0.5$\,s, bounding lead time above by $-2.4$\,s. The
target region is empty. We replace lead time by \emph{detection delay}.

\paragraph{Detection delay is dominated by acquisition.} Every detector we
evaluate reports an identical delay of $8.78$\,s. The cause is structural: the
first burst of the fault file is $4$ samples long and admits no window, so the
first evaluable window begins $8.78$\,s after $t_0$. Detection delay therefore
measures the data, not the model.

\paragraph{Event recall is constant.} With a single positive event, every
detector that fires at all scores $1.0$. The Clopper--Pearson interval is
$[0.025, 1.0]$.

\paragraph{False-alarm rate is unresolvable.} The normal test segment provides
$0.139$\,h of data time. Observing zero alarms gives a Poisson $95\%$ upper
bound of $21.6$ alarms/h; one alarm gives $34.2$/h. A target of
$\leq 1$ alarm/h would require roughly $3$\,h of alarm-free normal observation.
Treating ``$0.0$/h'' as a hard admission criterion therefore filters candidates
on noise. We report false-positive \emph{counts} and use them only for ranking.

\subsection{What remains}

After the above, the identifiable axes are the bootstrap confidence interval of
window $F_1$, the false-positive count, and model complexity. Because stride-1
windows overlap by $19/20$ samples, we resample at the burst level: all
confidence intervals and paired comparisons use a block bootstrap over bursts
($B=1000$), which is the smallest unit that is plausibly exchangeable.

\subsection{Two-tier baselines}

A zero-parameter rule that flags any window leaving the per-channel training
range attains $F_1 = 0.857$ under the reference protocol, against $0.748$ for
the LSTM-autoencoder \citep{malhotra2016lstm} distributed with the dataset.
Window $F_1$ is thus saturated near $0.85$, and a deep model scoring $0.80$ is
not ``good'' but worse than no model.

However, promoting this rule to \emph{the} baseline is also wrong, for a reason
specific to the data: the extrema of all three channels occur before row
$15{,}000$, so the test segment lies strictly inside the training range and the
rule \emph{cannot} produce a false positive. Its $\text{FPR}=0$ is an identity,
not a measurement. We therefore report two tiers---the conventional reference
model and the strict zero-parameter rule---and we verify, for every range-based
detector we report, that the training range is in fact exceedable on test data
(it is, for $2.6$--$3.8\%$ of windows).

\section{Session Confound and a Pre-declared Exclusion Rule}
\label{sec:confound}

\subsection{The current channel separates the sessions, not the classes}

The lag-1 autocorrelation of \textsf{AI2} is $0.927$ in the normal file with a
minimum of $0.906$, and at most $0.836$ in the fault file: the two classes are
perfectly separated by one scalar. Taken at face value this is the strongest
feature in the dataset. Three observations argue it is an artefact.

First, normal \textsf{AI2} is a pure aliased tone. A single sinusoid fits each
burst with median $R^2 = 0.996$ at an apparent $0.600$\,Hz, and
$\cos(2\pi\cdot0.6/10) = 0.9298$ reproduces the observed autocorrelation. This
identifies a $60$\,Hz mains carrier folded by a sampling clock in error by
$1.0\%$. Second, the statistic is a function of carrier \emph{frequency}, which
load cannot change, so its invariance across operating conditions carries no
information about whether it tracks machine health. Third, and decisively, the
share of AC power in the carrier band falls from $0.959$ (normal, range
$[0.947,0.973]$) to $0.104$ (fault, $[0.046,0.183]$) with no overlap. For a
running induction motor the mains fundamental cannot vanish; a change of this
magnitude indicates a different acquisition path.

We also note a file-level fingerprint independent of any channel: decimal
precision differs between the files (median $6$ digits with $0.0\%$ exceeding
seven in the normal file, versus median $8$ and $64.8\%$ in the fault file,
including float64 serialisation residue such as
\texttt{-44.107442000000006}). The two files were exported by different
pipelines. The confound is therefore present at the file level and is not
specific to any one variable.

\subsection{The rule}

Because a single fault session cannot falsify an instrumental explanation, we
decline to adjudicate it and instead pre-declare an exclusion rule:

\begin{quote}\small
\textbf{D1.} A feature that is invariant (coefficient of variation $<1\%$)
across all observed normal operating variation, yet differs between acquisition
sessions, is excluded from the final model, because a single session cannot
falsify an instrumental explanation.
\end{quote}

D1 is a procedural rule, not a physical verdict. It applies to the whole
\textsf{AI2} channel, including the amplitude features \textsf{AI2\_std} and
\textsf{AI2\_ptp}, which are functions of the same carrier. We regard partial
exclusion---banning the autocorrelation while retaining the amplitude---as
indefensible, and we report the cost of full exclusion in
\S\ref{sec:results}.

\section{Features and Detectors}
\label{sec:method}

\subsection{Feature sets}

All features are computed on gap-aware windows of $L=20$. We evaluate four
D1-compliant sets, each built to test one hypothesis about what distinguishes
the fault, plus two D1-violating sets retained only for sensitivity analysis.

\begin{itemize}\itemsep1pt
\item \textsf{S0\_AMP4}: $\{\sigma, \text{p2p}\}$ for \textsf{AI0}, \textsf{AI1}. Amplitude only.
\item \textsf{S1\_VIB}: \textsf{S0} $+$ lagged correlations at lags $1,2$ $+$
$\rho(\textsf{AI0},\textsf{AI1})$ $+$ amplitude ratio $+$ interaction.
\item \textsf{S2\_SHAPE}: amplitude $+$ kurtosis, crest factor $+$ ratio $+$ interaction.
\item \textsf{S3\_SPEC}: amplitude $+$ band power at $1.5$--$2.5$ and $3.0$--$4.0$\,Hz $+$ ratio $+$ interaction.
\end{itemize}

Two constructions deserve comment. The \emph{amplitude ratio}
$\sigma(\textsf{AI0})/\sigma(\textsf{AI1})$ is invariant to per-channel gain and
polarity, which matters because the dataset documents no sensor axis, mounting
angle, or unit; consequently no absolute vibration-severity standard can be
applied, and all thresholds must be relative to the data itself. The
\emph{interaction} term $-\rho\cdot\sigma(\textsf{AI0})$ encodes a finding from
\S\ref{sec:fp}: antiphase motion alone is not diagnostic because it also occurs
in normal operation; antiphase \emph{together with} increased amplitude is.

\subsection{Detectors}

All detectors are one-class and are fitted on normal training windows only:
the range rule; Mahalanobis distance and a one-sided variant that accumulates
only positive deviations; principal-component $T^2$ and squared prediction
error \citep{jackson1979control,macgregor1995statistical}; Isolation Forest
\citep{liu2008isolation}; Local Outlier Factor in novelty mode
\citep{breunig2000lof}; and a one-class SVM \citep{scholkopf2001estimating}.
Thresholds are the $q_{99.9}$ quantile of the score on the normal validation
segment, so no fault label enters threshold selection. Scores are converted to
$p$-values by split conformal prediction \citep{vovk2005algorithmic,
laxhammar2015inductive}, which requires only normal data; we use the
uniformity of normal $p$-values as a calibration diagnostic rather than
claiming a coverage guarantee, since exchangeability is violated by the
excursion of \S\ref{sec:fp}.

\section{Results}
\label{sec:results}

\subsection{Protocol}

Train on normal rows $[0,12000)$ ($6{,}021$ windows), select thresholds on
normal rows $[12000,15000)$ ($1{,}468$ windows), and test on normal rows
$[15000,20000)$ ($2{,}489$ windows) together with all $276$ fault windows.
Standardisation is fitted on the training segment only and is omitted for the
range rule and Isolation Forest, which are invariant to monotone per-feature
transforms. Confidence intervals are burst-level block bootstraps with
$B=1000$.

\begin{table}[t]
\centering\small
\caption{D1-compliant detectors, ranked by the lower bound of the bootstrap
$F_1$ interval. FP is the count of false positives among $2489$ normal test
windows; $|\theta|$ is the number of learned parameters.}
\label{tab:main}
\begin{tabular}{llcccc}
\toprule
Features & Detector & $F_1$ & 95\% CI & FP & $|\theta|$ \\
\midrule
\textsf{S1\_VIB} & Range rule & \textbf{0.928} & [.852,.984] & \textbf{0} & \textbf{0} \\
\textsf{S3\_SPEC} & PCA-SPE & 0.917 & [.849,.971] & 10 & 55 \\
\textsf{S3\_SPEC} & Mahalanobis & 0.905 & [.829,.970] & 11 & 65 \\
\textsf{S2\_SHAPE} & OCSVM & 0.903 & [.817,.976] & 17 & 95 \\
\textsf{S2\_SHAPE} & LOF & 0.896 & [.808,.973] & 4 & 0 \\
\textsf{S2\_SHAPE} & Mahalanobis$^{1S}$ & 0.888 & [.794,.971] & 22 & 65 \\
\textsf{S2\_SHAPE} & Range rule & 0.877 & [.786,.955] & 6 & 0 \\
\textsf{S1\_VIB} & PCA-$T^2$ & 0.863 & [.791,.927] & 60 & 72 \\
\textsf{S0\_AMP4} & OCSVM & 0.866 & [.773,.949] & 7 & 72 \\
\bottomrule
\end{tabular}
\end{table}

\subsection{Main result}

Table~\ref{tab:main} gives the D1-compliant leaderboard. The best detector is a
zero-parameter range rule on \textsf{S1\_VIB}, at $F_1 = 0.928$ with perfect
precision and no false positives. Against the strict baseline of
\S\ref{sec:metrics} ($F_1 = 0.818$) the paired block bootstrap gives
$\Delta F_1 \in [+0.008, +0.202]$ with $P(\Delta \leq 0) = 0.019$; against a
Mahalanobis detector on amplitude features only ($0.828$) it gives
$[+0.055, +0.159]$, $P(\Delta\leq0)=0.000$.

We verify that this result is not the identity of \S\ref{sec:metrics}: in the
\textsf{S1\_VIB} feature space, $95$ of $2489$ normal test windows ($3.82\%$)
do exceed the training range on at least one coordinate, and are nonetheless
scored below threshold.

The second-ranked detector, PCA-SPE on spectral features, is
\emph{not distinguishable} from the first: $\Delta F_1 \in [-0.022, +0.043]$,
$P(\Delta\leq0)=0.265$. We consider this the most useful single number in the
study. Rather than break the tie on point estimates, we assign the two models
different roles: the zero-parameter rule is the \emph{reference} that any more
complex model must beat significantly, and the PCA model is the
\emph{deployment} candidate, because its residual decomposition yields
contribution plots that attribute an alarm to individual variables---the
standard diagnostic artefact in multivariate process control.

\subsection{Cost of the exclusion rule}

\begin{table}[t]
\centering\small
\caption{Sensitivity to D1. Excluding the current channel costs
$\Delta F_1 = 0.072$. Two current-derived features alone separate the classes
perfectly, which we read as a measure of the confound rather than of
performance.}
\label{tab:d1}
\begin{tabular}{llcc}
\toprule
Features & Detector & $F_1$ & D1 \\
\midrule
\textsf{AI2} only (2 feat.) & PCA-SPE & 1.000 & violates \\
\textsf{AI2} only (2 feat.) & Mahalanobis & 1.000 & violates \\
amp.\ $+$ \textsf{AI2} & LOF & 0.984 & violates \\
amp.\ $+$ \textsf{AI2} & OCSVM & 0.964 & violates \\
\midrule
\textsf{S1\_VIB} (adopted) & Range rule & 0.928 & complies \\
\bottomrule
\end{tabular}
\end{table}

Table~\ref{tab:d1} quantifies what D1 costs. Two current-derived features
suffice for perfect separation. We report this prominently precisely because it
is the exclusion's justification: given independent evidence that the current
acquisition path differed between sessions, perfect separation measures the
confound, not the detector. The cost of declining it is $\Delta F_1 = 0.072$.

\subsection{Where the false alarms are}
\label{sec:fp}

\begin{table}[t]
\centering\small
\caption{False alarms are not uniformly distributed over the normal test
segment. Rows are $1000$-sample bins; statistics are bin means. The detector is
the worst-performing D1-compliant candidate ($191$ false positives), chosen to
make the concentration visible.}
\label{tab:fp}
\begin{tabular}{lrrrrr}
\toprule
Rows & $n$ & FP & rate & $\sigma_{\textsf{AI1}}$ & $\rho_{01}$ \\
\midrule
15--16k & 501 & 13 & 2.6\% & 0.089 & $+0.197$ \\
\textbf{16--17k} & 499 & \textbf{64} & \textbf{12.8\%} & \textbf{0.043} & $\mathbf{-0.198}$ \\
\textbf{17--18k} & 504 & \textbf{56} & \textbf{11.1\%} & \textbf{0.042} & $\mathbf{-0.248}$ \\
\textbf{18--19k} & 502 & \textbf{54} & \textbf{10.8\%} & 0.058 & $-0.093$ \\
19--20k & 483 & 4 & 0.8\% & 0.103 & $+0.316$ \\
\bottomrule
\end{tabular}
\end{table}

Table~\ref{tab:fp} shows that $91.1\%$ of false positives fall in rows
$16{,}000$--$19{,}000$, roughly $300$\,s. In that interval vibration amplitude
drops by $53\%$ and the upper--lower correlation reverses sign, yet the label
remains normal; by row $19{,}000$ both statistics return to training levels
($0.103$ and $+0.316$ against $0.115$ and $+0.306$). This is a reversible
second operating mode, not a regime change, and we were able to confirm that it
co-varies with a load proxy.

Two implications follow. First, sign reversal of inter-sensor correlation is
\emph{not} a fault signature: it occurs in labelled-normal data. What
distinguishes the fault is reversal \emph{accompanied by rising} amplitude,
whereas the normal excursion reverses while amplitude \emph{falls}---which is
why we include the interaction feature of \S\ref{sec:method}. Second, and less
comfortably, the zero false positives of our adopted detector are a consequence
of direction rather than robustness: normal operating variation pushes
amplitude \emph{inside} the training support while the fault pushes it outside.
A normal mode that raised amplitude---a heavier blanking job, a higher pressure
setting---would produce false alarms. No such mode is present in these data,
and we state this as a limitation rather than discovering it later.

\section{Discussion}

\paragraph{Ordering of decisions.} Across every comparison we ran, the choice
of windowing ($\Delta F_1 = 0.039$), the choice of whether to admit the
confounded channel ($0.072$), and the choice of baseline (the difference
between claiming $+0.18$ and claiming $+0.11$ over the reference model) each
moved the headline number by more than the spread across all eight detectors on
a fixed feature set. Model selection was the least consequential decision we
made. We suspect this ordering is common in small industrial benchmarks and is
under-reported because the preprocessing decisions are rarely ablated.

\paragraph{On zero-parameter winners.} That a range rule tops
Table~\ref{tab:main} is not an argument against learning; it is a statement
about the identifiability of this benchmark. With $276$ positive windows drawn
from one event, the effective positive sample size is one, and the bootstrap
intervals in Table~\ref{tab:main} overlap almost completely. The honest reading
is that this dataset can rank preprocessing choices but cannot rank detectors.

\paragraph{Calibration versus coverage.} We deliberately do not claim a
conformal coverage guarantee. Exchangeability fails across the excursion of
\S\ref{sec:fp}, and in feature sets that include correlation statistics we
measured empirical false-positive rates exceeding the nominal $\alpha$ by up to
$84\times$. We report the Kolmogorov--Smirnov distance of the normal
$p$-value distribution from uniform as a diagnostic instead; it is smallest
($0.046$--$0.105$) exactly for the amplitude and spectral feature sets that are
insensitive to the excursion.

\paragraph{Limitations.} The study concerns one dataset with one fault event,
and all statements about the fault mechanism are necessarily weak. We make no
claim about which failure mode produced the recording; the available bandwidth
($\leq 5$\,Hz) excludes the diagnostic frequencies that would be required. Our
evaluation protocol also inherits the dataset's single-session structure: a
model that passes D1 may still be exploiting session differences that we failed
to identify.

\section{Conclusion}

On a representative industrial predictive-maintenance benchmark, we find that
three evaluation decisions---gap-aware windowing, baseline selection, and the
treatment of a session-confounded channel---each affect the reported $F_1$ more
than the choice among eight standard one-class detectors, and that three
commonly reported operational metrics are unidentifiable when the positive
class is a single event. We recommend that work on such benchmarks report the
windowing survival rate, verify that range-based baselines can in principle
produce false positives, state the effective number of positive events
alongside any interval estimate, and quantify the cost of excluding any channel
suspected of acquisition confounding. Applying this protocol, a vibration-only
detector with no learned parameters attains $F_1 = 0.928$ with zero false
positives, and $91.1\%$ of all false alarms in the study localise to a single
reversible load excursion that the benchmark labels normal---which we regard as
the finding the data actually supports.

\paragraph{Reproducibility.} The full pipeline runs end to end from a single
entry point and emits the comparison table, the false-alarm decomposition, and
per-window predictions. All randomised components are seeded; the range rule,
Mahalanobis, and PCA detectors are deterministic. Repeated runs are
byte-identical.

\begin{thebibliography}{99}\small
\setlength{\itemsep}{1pt}

\bibitem[Breunig et al.(2000)]{breunig2000lof}
M.~M. Breunig, H.-P. Kriegel, R.~T. Ng, and J.~Sander.
\newblock LOF: Identifying density-based local outliers.
\newblock In \emph{SIGMOD}, 2000.

\bibitem[Garg et al.(2021)]{garg2021evaluation}
A.~Garg, W.~Zhang, J.~Samaran, R.~Savitha, and C.-S. Foo.
\newblock An evaluation of anomaly detection and diagnosis in multivariate time
  series.
\newblock \emph{IEEE Trans. Neural Netw. Learn. Syst.}, 2021.

\bibitem[Hundman et al.(2018)]{hundman2018detecting}
K.~Hundman, V.~Constantinou, C.~Laporte, I.~Colwell, and T.~Soderstrom.
\newblock Detecting spacecraft anomalies using LSTMs and nonparametric dynamic
  thresholding.
\newblock In \emph{KDD}, 2018.

\bibitem[Jackson \& Mudholkar(1979)]{jackson1979control}
J.~E. Jackson and G.~S. Mudholkar.
\newblock Control procedures for residuals associated with principal component
  analysis.
\newblock \emph{Technometrics}, 21(3):341--349, 1979.

\bibitem[Kim et al.(2022)]{kim2022rigorous}
S.~Kim, K.~Choi, H.-S. Choi, B.~Lee, and S.~Yoon.
\newblock Towards a rigorous evaluation of time-series anomaly detection.
\newblock In \emph{AAAI}, 2022.

\bibitem[Laxhammar \& Falkman(2015)]{laxhammar2015inductive}
R.~Laxhammar and G.~Falkman.
\newblock Inductive conformal anomaly detection for sequential detection of
  anomalous sub-trajectories.
\newblock \emph{Ann. Math. Artif. Intell.}, 74:67--94, 2015.

\bibitem[Liu et al.(2008)]{liu2008isolation}
F.~T. Liu, K.~M. Ting, and Z.-H. Zhou.
\newblock Isolation forest.
\newblock In \emph{ICDM}, 2008.

\bibitem[Ljung \& Box(1978)]{ljung1978measure}
G.~M. Ljung and G.~E.~P. Box.
\newblock On a measure of lack of fit in time series models.
\newblock \emph{Biometrika}, 65(2):297--303, 1978.

\bibitem[MacGregor \& Kourti(1995)]{macgregor1995statistical}
J.~F. MacGregor and T.~Kourti.
\newblock Statistical process control of multivariate processes.
\newblock \emph{Control Engineering Practice}, 3(3):403--414, 1995.

\bibitem[Malhotra et al.(2016)]{malhotra2016lstm}
P.~Malhotra, A.~Ramakrishnan, G.~Anand, L.~Vig, P.~Agarwal, and G.~Shroff.
\newblock LSTM-based encoder-decoder for multi-sensor anomaly detection.
\newblock \emph{arXiv:1607.00148}, 2016.

\bibitem[Sch\"olkopf et al.(2001)]{scholkopf2001estimating}
B.~Sch\"olkopf, J.~C. Platt, J.~Shawe-Taylor, A.~J. Smola, and R.~C.
  Williamson.
\newblock Estimating the support of a high-dimensional distribution.
\newblock \emph{Neural Computation}, 13(7):1443--1471, 2001.

\bibitem[Vovk et al.(2005)]{vovk2005algorithmic}
V.~Vovk, A.~Gammerman, and G.~Shafer.
\newblock \emph{Algorithmic Learning in a Random World}.
\newblock Springer, 2005.

\bibitem[Wu \& Keogh(2021)]{wu2021flawed}
R.~Wu and E.~Keogh.
\newblock Current time series anomaly detection benchmarks are flawed and are
  creating the illusion of progress.
\newblock \emph{IEEE Trans. Knowl. Data Eng.}, 2021.

\end{thebibliography}

\end{document}
